% 树的性质示意：(a) 两条不同u-v路径在x分开、在y会合形成环；(b) 最长路径端点的另一个邻居导致矛盾；
% (c) 同一棵树（边{1,2},{1,3},{3,4},{3,5}）分别以1和4为根。
\begin{tikzpicture}[>=Stealth,
  v/.style={circle,draw=FigureStroke,fill=FigureBlueFill,line width=.8pt,minimum size=6mm,inner sep=0pt},
  vr/.style={v,draw=FigureGreen,fill=FigureGreenFill},
  e/.style={draw=FigureStroke,line width=.8pt},
  pb/.style={draw=FigureBlue,line width=1.4pt},
  pg/.style={draw=FigureGreen,line width=1.4pt,dash pattern=on 3pt off 1.5pt},
  bad/.style={draw=FigureRed,dashed,line width=.9pt},
  note/.style={font=\FigureNoteFont\fontsize{7.5}{9}\selectfont,text=FigureMuted,align=center}]
% panel a
\begin{scope}
  \fill[FigureYellowFill] (1,0) .. controls (1.6,1.0) and (2.4,1.0) .. (3,0)
                          .. controls (2.4,-1.0) and (1.6,-1.0) .. (1,0);
  \node[v] (u) at (0,0) {$u$}; \node[v] (x) at (1,0) {$x$};
  \node[v] (y) at (3,0) {$y$}; \node[v] (w) at (4,0) {$v$};
  \draw[e] (u)--(x); \draw[e] (y)--(w);
  \draw[pb] (x) .. controls (1.6,1.0) and (2.4,1.0) .. (y);
  \draw[pg] (x) .. controls (1.6,-1.0) and (2.4,-1.0) .. (y);
  \node[text=FigureBlue] at (2,1.0) {$P$}; \node[text=FigureGreen] at (2,-1.0) {$Q$};
  \node[note] at (2,0) {构成环};
  \node[align=center] at (2,-1.6) {(a) 两条路径在$x$分开、在$y$会合};
\end{scope}
% panel b
\begin{scope}[xshift=6.2cm]
  \node[v] (0) at (0,0) {$v_0$}; \node[v] (1) at (1,0) {$v_1$};
  \node[v] (2) at (2,0) {$v_2$}; \node[v] (3) at (3,0) {$v_3$};
  \draw[pb] (0)--(1)--(2)--(3);
  \node[v,draw=FigureRed,fill=FigureRedFill] (z) at (-.7,-.8) {$z$};
  \draw[bad] (0)--(z);
  \draw[bad] (0) to[bend left=45] (2);
  \node[note,anchor=west] at (-.3,-.85) {$z$不在路径上：路径可加长};
  \node[note] at (1,1.0) {邻居在路径上：成环};
  \node[align=center] at (1.5,-1.6) {(b) 最长路径的端点是叶};
\end{scope}
% panel c
\begin{scope}[yshift=-2.6cm,xshift=1.2cm]
  \node[vr] (1) at (0,0) {$1$}; \node[v] (2) at (-.6,-.75) {$2$}; \node[v] (3) at (.6,-.75) {$3$};
  \node[v] (4) at (.1,-1.5) {$4$}; \node[v] (5) at (1.1,-1.5) {$5$};
  \draw[e] (1)--(2) (1)--(3) (3)--(4) (3)--(5);
  \node at (.25,-2.75) {以$1$为根：叶$2,4,5$};
\end{scope}
\begin{scope}[yshift=-2.6cm,xshift=4.6cm]
  \node[vr] (4) at (0,0) {$4$}; \node[v] (3) at (0,-.75) {$3$};
  \node[v] (1) at (-.6,-1.5) {$1$}; \node[v] (5) at (.6,-1.5) {$5$}; \node[v] (2) at (-.6,-2.25) {$2$};
  \draw[e] (4)--(3) (3)--(1) (3)--(5) (1)--(2);
  \node at (0,-2.75) {以$4$为根：叶$2,5$};
\end{scope}
\node[note,anchor=west,align=left] at (6.3,-3.6) {两图边集相同，\\选根只改变父子关系；\\度为$1$的根不算叶};
\node[align=center] at (3.2,-5.75) {(c) 同一棵树选不同的根};
\end{tikzpicture}
